\documentclass[11pt]{article}

% --------------------------------------------------
% Packages
% --------------------------------------------------
\usepackage[margin=1in]{geometry}
\usepackage{amsmath, amssymb, amsthm}
\usepackage{mathtools}
\usepackage{graphicx}
\usepackage{float}
\usepackage{enumitem}
\usepackage{hyperref}
\usepackage{xcolor}
\usepackage{tikz}
\usetikzlibrary{arrows.meta, positioning, calc}

% Pseudocode and algorithms
\usepackage{algorithm}
\usepackage{algpseudocode}

% Better fonts (optional)
\usepackage{lmodern}

% --------------------------------------------------
% Hyperref setup
% --------------------------------------------------
\hypersetup{
    colorlinks=true,
    linkcolor=blue,
    citecolor=blue,
    urlcolor=blue
}

% --------------------------------------------------
% Theorem-like environments
% --------------------------------------------------
\theoremstyle{definition}
\newtheorem{definition}{Definition}[section]
\newtheorem{theorem}{Theorem}[section]
\newtheorem{claim}{Claim}[section]
\newtheorem{remark}{Remark}[section]
\newtheorem{recall}{Recall}[section]

\theoremstyle{proof}
% Proof environment is built-in via amsthm

% --------------------------------------------------
% Custom commands
% --------------------------------------------------
\newcommand{\Gen}{\mathsf{Gen}}
\newcommand{\Enc}{\mathsf{Enc}}
\newcommand{\Dec}{\mathsf{Dec}}
\newcommand{\Sig}{\mathsf{Sig}}
\newcommand{\Ver}{\mathsf{Ver}}
\newcommand{\negl}{\mathsf{negl}}

% --------------------------------------------------
% Document
% --------------------------------------------------
\begin{document}

\title{Public Key Crypto notes - 08}
\author{By contributors}
\date{March 30. 2026}
\maketitle

% --------------------------------------------------
\section{Oblivious transfer recap}
% --------------------------------------------------

$A$ sends $m_0, m_1$, $B$ chooses $b \in_R \{0, 1\}$ and obtains $m_b$. RSA-based protocol.

\section{Protocol 2 -- Bellari-Micali}

\begin{itemize}
	\item Diffie-Hellman-based, interactive protocol.
	\item Let $p, g, q$ be as in ElGamal.
	\item $A$ chooses $c \in_R <g>$, where $<g> = \{g^i: i = 0, \dots, q-1\}$ and sends $c$ to $B$.
	\item $B$ chooses $k \in_R \mathbb{Z}_q, b \in_R \{0, 1\}, y_b = g^k, y_{1-b} = \frac{c}{g^k}$, and sends $y_0, y_1$ to $A$.
	\item $A$ checks $y_0 \cdot y_1 \stackrel{?}{=} c$ and chooses two random exponents $r_0, r_1 \in_R \mathbb{Z}_q$
	\begin{itemize}
		\item $c_0 = (g^{r_0}, H(y_0^{r_0}) \oplus m_0)$
		\item $c_1 = (g^{r_1}, H(y_1^{r_1}) \oplus m_1)$
		\item $A$ sends $c_0, c_1$ to $B$
	\end{itemize}
	\item $B$ calculates $c_b = (v_0, v_1), m_b = H(v_0^k) \oplus v_1$
\end{itemize}

\section{Application of oblivious transfers}

\begin{itemize}
	\item Secure multiparty computations (MPC)
	\begin{itemize}
		\item There are two parties $A, B$. Both of them have a secret input $x_A, x_B$. There is a function $f(x, x')$. $A$ and $B$ want to evaluate $f(x_A, x_B)$ without them learning the other's secret.
		\item Example: Dating protocol
		\begin{itemize}
			\item Two parties: $A, B$
			\item Let $d_A = \begin{cases}
				1 \text{ if } A \text{ wants to date with B} \\
				0 \text{ otherwise}
			\end{cases}$ and, $d_B = \begin{cases}
			1 \text{ if } B \text{ wants to date with A} \\
			0 \text{ otherwise}
			\end{cases}$
			\item The function to evaluate this: $\text{AND}(d_A, d_B)$
		\end{itemize}
		\item Evaluating AND with oblivious transfer:
		\begin{itemize}
			\item If $d_A = 1$, let $m_0 = 0$ (meaning that $B$ chose $0$), $m_1 = 1$
			\item If $d_A = 0$, let $m_0 = 0, m_1 = 0$
			\item Then, $A$ sends $m_0, m_1$ to $B$ by oblivious transfer
		\end{itemize}
	\end{itemize}
\end{itemize}

\section{Group signatures}

\begin{itemize}
	\item Group/organization
	\item A Group Manager can add users to a group
	\item The group $(S_0, S_1, \dots S_n)$ jointly issues signature $(m, \sigma)$ on behalf of the group itself
	\item The exact individual who actually made the signature is masked
	\item Requirements:
	\begin{itemize}
		\item $\sigma$ cannot be linked to any $S_i$
		\item The Group Manager $M$ can resolved $\sigma$ (i.e. can learn the actual signer $S_i$)
	\end{itemize}
	\item Algorithms:
	\begin{itemize}
		\item Setup: $M$ sets the parameters up
		\item Joining member $S_i$ is added by $M$ (can even happen dynamically)
		\item Signing: $S_i$ signs $m$ and outputs $\sigma$
		\item Verification: Third-party user $U$ verifies by $\Ver(m, \sigma)$
		\item Open: $M$ links the signature $\sigma$ to $S_i$
	\end{itemize}
\end{itemize}

\subsection{Examples of group signature protocols}

\subsubsection{Chaum's group signature, '91}

\begin{itemize}
	\item Setup/Join:
	\begin{itemize}
		\item $S_i$ generates ElGamal keys: $pk_i = (p, g, q, g^{x_i}), sk_i = x_i$
		\item $S_i$ gives $pk_i$ to $M$
	\end{itemize}
	\item Sign:
	\begin{itemize}
		\item At period $t$, $M$ chooses $r_{t,i} \in_R \mathbb{Z}_q$, and sends it to $S_i$
		\item $M$ publishes $P_t = \{g^{r_{t,i} \cdot x_i}, i = 1, 2, \dots\}$ (all blinded public keys at period $t$)
		\item $S_i$ uses $r_{t, i} \cdot x_i$ as a secret key for the ElGamal signature, $pk = g^{r_{t, i} \cdot x_i}$.
	\end{itemize}
	\item Verify:
	\begin{itemize}
		\item $U$ checks: $pk \stackrel{?}{\in} P_t$ and verifies the ElGamal signature.
	\end{itemize}
	\item Open:
	\begin{itemize}
		\item $M$ checks for which $i: pk = (g^{x_i})^{r_{t, i}}$
	\end{itemize}
\end{itemize}

\begin{remark}
	At period $t$, only \emph{one} signer can make a signature. Hence, the whole signature $(m, \sigma, P_t)$ is linear in the size of the group.
\end{remark}

Before protocol 2, we need preparation with proof of knowledge protocols.

\section{Proof of knowledge protocols}

\subsection{Schnorr's PoK protocol}

\begin{itemize}
	\item Schnorr $\text{PoK}[x: g^x = h](m)$	
	\item Prover chooses $k \in_R \mathbb{Z}_q$, sends $I = g^k$ to the Verifier.
	\item Verifier chooses $r \in_R \mathbb{Z}_q \text{ or } r = H(I) \text{ or } r = H(I || m)$, and sends $R$ to the Prover
	\item Prover calculates $s = rx + k$, and sends $s$ to the Verifier
	\item The Verifier checks whether $r \stackrel{?}{=} H(g^s \cdot h^{-r} || m)$, where $g^s \cdot h^{-r} = I$
\end{itemize}


\begin{remark}
	Success probability without knowing $x$ is $\approx \frac{1}{q}$.
\end{remark}

\subsection{PoK of LogLog}

Goal: Prover knows $x$ such that $y = g^{(a^x)}$ with given $g, y, a$.

\begin{itemize}
	\item Prover chooses $k \in_R \{0, \dots, 2^\lambda\}, t = g^{a^k}$ and $t$ is sent to the Verifier
	\item The Verifier chooses $r \in_R \{0, 1\}$ and sends it to the Prover
	\item The Prover calculates: $s = \begin{cases}
		k, \text{ if } r = 0 \\
		k-x, \text{ if } r = 1
	\end{cases}$ and sends $s$ to the Verifier
	\item Verifier performs the following: $t = \begin{cases}
		g^{(a^s)}, \text{ if } r = 0 \\
		y^{(a^s)}, \text{ if } r = 1
	\end{cases}$
	\item Correctness:
	\begin{itemize}
		\item $r = 0$ \checkmark
		\item $r = 1 \rightarrow y^{(a^s)} = y^{(a^{k-x})} = [g^{(a^x)}]^{(a^{k-x})} = g^{(a^x) \cdot (a^k) \cdot (a^{-x})} = g^{(a^k)} = t$
	\end{itemize}
\end{itemize}

\subsubsection{Remarks}

\begin{remark}
	If $r = 0$, the Verifier checks if $t$ has the right form. If $r = 1$, then the Verifier checks if $P$ knows $x$, provided $t$ has the right form.
\end{remark}

\begin{remark}
	Success probability is $\frac{1}{2}$.
\end{remark}

\subsection{Non-interactive PoK protocol}

PoK$[x: y = g^{(a^x)}](m = r, s_1, \dots, s_l) \in \{0, B^l \cdot \mathbb{Z}_q^l\}$ such that $r = H(m || y || g || a || t_1 || \dots || t_l)$ with $t_i = \begin{cases}
	g^{(a^{s_i})}, \text{ if } r_i = 0 \\
	y^{(a^{s_i})}, \text{ if } r_i = 1.
\end{cases}$

\noindent Construction:

\begin{itemize}
	\item $t_i = g^{(a^{k_i})}$ with random $k_i$ and $s_i = \begin{cases}
		k_i, \text{ if } r_i = 0 \\
		k_{i-x}, \text{ if } r_i = 1
	\end{cases}$.
\end{itemize}

\end{document}